\section{Related work}
\label{sec:related}

Prior work on football roles has generally required data that is not public. Bialkowski et
al. \citep{bialkowski2014large} recovered formations and positional roles from tracking
coordinates. Decroos and Davis \citep{decroos2019player} learned player vectors from event
streams, Pappalardo et al. \citep{pappalardo2019playerank} built a performance rating from
event data across several leagues, and Aalbers and Van Haaren
\citep{aalbers2019distinguishing} distinguished playing styles from match events. All
depend on proprietary feeds. Where public season aggregates are used instead, the standard
output is a taxonomy. \citet{elfoukahi2026unveiling} apply $k$-means to twelve indicators
from the same source used here and report four midfielder archetypes, and
\citet{seri2024partial} fit partial membership models to a season of Serie A counts. Two
systematic reviews catalogue the genre, \citet{plakias2023identifying} for players and
\citet{gu2025playing} for teams; the first synthesises twenty-six styles from twelve
studies and finds no shared framework behind them. The pattern is not specific to
football. \citet{alagappan2012from} reported thirteen positions in basketball from a
topological summary of box score rates, and \citet{kalman2020nba} nine from a Gaussian
mixture. In none of this work is it asked whether the data contain groups at all.
\citet{fernandez2024reporting}, reviewing 278 sports science articles that used
clustering, found that 55 percent did not state the criterion by which the number of
clusters was chosen and 27 percent did not state the clustering method, which puts the
omission in context: the question of existence has not been answered badly so much as not
asked.

The closest methodological relative of the present paper is the work of Akhanli and
Hennig. Their aggregated validity index \citep{akhanli2020comparing} combines several
criteria, each calibrated by comparing its value on a candidate partition against its
distribution over random partitions of the same data, and applied to a season of player
statistics from eight European leagues \citep{akhanli2023clustering} it returned a coarse
partition into five groups, ordered from centre backs through full backs, midfielders and
attacking midfielders to forwards, together with a fine partition into about one hundred
and fifty small clusters intended for similarity search. The coarse partition runs along
the same territorial axis reported here. The difference is in what the calibration is
against. A random partition of the observed points asks how good a clustering is relative
to other clusterings of the same data, which is the right question for choosing among
methods and numbers of clusters. Clusterless data matched to the observed sample asks
whether any clustering of these points is better than a clustering of points that have no
structure, which is the question of existence. The two are complementary, and the first
does not answer the second.

The methodological position taken here has a longer history than the football
literature, and the paper should be precise about what is and is not new. The problem of
choosing the number of clusters has been studied systematically since
\citet{milligan1985examination} compared thirty stopping rules on data with known
structure, and the silhouette \citep{rousseeuw1987silhouettes}, the gap statistic
\citep{tibshirani2001estimating} and the mixture information criterion all descend from
that programme. The gap statistic is the direct ancestor of the calibration used here: it
compares the observed within-cluster dispersion against the same quantity on reference
data drawn uniformly over the principal box, which is one of the three nulls in this
paper. What this paper adds to it is not a new statistic but a harder reference. A
uniform box is easy to beat, a Gaussian matched on covariance is harder, and a copula that
keeps every real marginal is harder still, and the paper reports the observed separation
against all three so that a reader can see which reference the claim survives.

Stability-based validation is the second tradition the paper engages, and it engages it
critically. \citet{benhur2002stability} proposed choosing the number of clusters by the
stability of partitions under subsampling, and the bootstrap adjusted Rand index reported
by most role taxonomies is a descendant of that idea. \citet{vonluxburg2010stability}
reviewed the theory and reached a conclusion that the applied literature has largely not
absorbed: for algorithms such as $k$-means, stability is governed by the geometry of the
objective rather than by the presence of clusters, so a partition can be perfectly stable
on data that contain no groups at all. Section~\ref{sec:results:structure} measures
exactly that on this data and finds it. The point is not that stability is useless, since
an unstable partition is certainly not a taxonomy, but that a stable one is not certainly
a taxonomy either, and the genre has been treating the second as if it followed from the
first.

Clusterability as a question distinct from clustering has its own literature.
\citet{ackerman2009clusterability} surveyed formal notions of it and found them pairwise
inconsistent, so that the same data can be well clustered under one definition and not
under another. \citet{adolfsson2019cluster} compared the practical tests and recommended
those built on Hartigan's dip statistic \citep{hartigan1985dip} applied to a low
dimensional projection, which is one of the instruments used here. \citet{hennig2015true}
argued that there is no true clustering independent of the purpose the clustering is meant
to serve, and that the honest question is whether a partition has the properties the
analyst needs rather than whether it is correct. This paper takes that position: the
property it asks for is that a partition be distinguishable from one imposed on a
continuum, and it reports which partitions have it.

A different tradition describes athletes as mixtures rather than as members. Archetypal
analysis \citep{eugster2012performance} represents each player as a convex combination of
a few extreme profiles, and partial membership models \citep{seri2024partial} let a player
belong to several clusters at once. Both are natural descriptions of a continuum, and both
presuppose one, since neither asks whether the data would have supported discrete groups
instead. The present paper supplies the test that decides between the two descriptions,
and its answer, two modes with a populated middle, is one that a mixture representation
expresses well and a taxonomy does not.

The role taxonomy genre has proceeded as though the cluster validity literature did not
exist. It reports partitions without asking whether the data are clusterable, and
validates them with criteria that, as Section~\ref{sec:results:structure} shows, a
clusterless continuum satisfies equally well. The contribution here is not a new
clusterability test but the application of existing ones to a field that has skipped
them, and the finding that they return a real but far coarser structure than the field has
been reporting. To our knowledge no published study in any sport has calibrated cluster
quality against data simulated to contain no clusters, measured bootstrap stability on
such data, or applied a clusterability test to player statistics, and the searches
conducted for this paper found none.
